\documentclass[12pt]{article}
\usepackage[T1]{fontenc}
\usepackage[utf8]{inputenc}
\usepackage{lmodern}
\usepackage{textcomp}
\usepackage{enumitem}
\usepackage{geometry}
\usepackage{graphicx}
\usepackage{amsmath, amssymb}
\geometry{margin=1in}
\begin{document}
\begin{center}
Constitutional Law - Graduate Level Assessment 11 \
Total Marks: 40 \
Answer all questions.
\end{center}

\section*{Multiple Choice (10 marks)}
\begin{enumerate}[label=\arabic*.]
\item Why is it important to separate the concept of punishment from its justification?
\begin{enumerate}[label=(\alph*)]
    \item Because the concept of punishment has evolved over time.
    \item Because punishment can be justified in multiple ways.
    \item Because any definition of punishment should be value-neutral.
    \item Because the practice of punishment is separate from its justification.
    \item Because the justification of punishment varies across cultures.
\end{enumerate}
\item In which instance would a state, under the enabling clause of the Fourteenth Amendment, be most able to regulate?
\begin{enumerate}[label=(\alph*)]
    \item A federal official from discriminating against a person based on race.
    \item A federal official from discriminating against a person based on gender.
    \item A federal official from discriminating against a person based on nationality.
    \item A private company from discriminating against a person based on nationality.
    \item A private individual from discriminating against a person based on race.
\end{enumerate}
\item The \_\_\_\_\_\_\_\_ School of jurisprudence postulates that the law is based on what is "correct."
\begin{enumerate}[label=(\alph*)]
    \item Legal Pragmatism
    \item Legal Formalism
    \item Comparative
    \item Analytical
    \item Sociological
\end{enumerate}
\item Which of the following will not terminate a contract by operation of law?
\begin{enumerate}[label=(\alph*)]
    \item Change in law that makes contract impossible
    \item Bankruptcy of one party
    \item Expiration of the contract term
    \item Partial performance
    \item Supervening illegality
\end{enumerate}
\item Catharine MacKinnon argues that since men dominate women, the question is ultimately one of power. Which proposition below is the most inconsistent with this argument?
\begin{enumerate}[label=(\alph*)]
    \item The power dynamic is determined by individual strength and not gender.
    \item The idea of 'woman' needs to be redefined.
    \item Women inherently possess more power than men.
    \item Gender-based inequality is a myth.
    \item Men and women share equal power in society.
\end{enumerate}
\end{enumerate}

\section*{True / False (10 marks)}
\textit{Answer True or False}
\begin{enumerate}[label=\arabic*.]
\setcounter{enumi}{5}
\item True or False: The idea of 'natural law' first emerged in Greek thinking, influencing later legal and philosophical frameworks.
\item True or False: Savigny's notion of Volksgeist suggests that law is primarily a tool used by the ruling class for societal oppression.
\item True or False: An 'armed attack' under Article 51 of the UN Charter encompasses all high-intensity instances of armed force.
\item True or False: Maine's aphorism that 'the movement of progressive societies has hitherto been a movement from Status to Contract' is often misinterpreted as a prediction.
\item True or False: Llewellyn's distinction between grand and formal styles of legal reasoning is misleading when it oversimplifies judicial roles.
\end{enumerate}

\section*{Long Form Response (20 marks)}
\textit{Answer each question with a clear and complete explanation.}
\begin{enumerate}[label=\arabic*.]
\setcounter{enumi}{10}
\item Discuss the Doctrine of Worthier Title and analyze its implications on property law, particularly focusing on how it transforms an invalid remainder in heirs into a reversion for the grantor.
\item Discuss the implications of Maine's aphorism that 'the movement of progressive societies has hitherto been a movement from Status to Contract' and analyze its common misinterpretation as a prediction.
\end{enumerate}

\vfill
\noindent\textit{End of Paper}
\end{document}